\scanpage{102}
\begin{center}- 10 -\end{center}
\margin{n° 285}

topologiques $(G_\alpha)$ et la famille d'applications $(j_\alpha\circ g_\alpha)$, en désignant par $g_\alpha$ l'application canonique de $H$ dans $H_\alpha$. Comme $j_\alpha\circ g_\alpha=f_\alpha\circ j$, le même critère de transitivité prouve que $j$ est un isomorphisme de $H$ sur le sous-groupe $j(H)$. En outre, $j(H)$ est l'intersection, dans le groupe produit $F=\prod_{\alpha\in I}G_\alpha$, de $G$ et du sous-groupe $\prod_{\alpha\in I}H_\alpha$ d'où la dernière assertion du corollaire.

Lorsque les hypothèses du cor.2 sont satisfaites, on dit que $(H_\alpha)$ est un \uline{système projectif de sous-groupes} correspondant au système projectif de groupes $(G_\alpha)$, et on identifie d'ordinaire $H$ et $j(H)$ au moyen de l'injection $j$.

Soit $K$ un sous-groupe de $G=\varprojlim G_\alpha$, et posons $K_\alpha=f_\alpha(K)$ ; il est immédiat que $(K_\alpha)$ est un système projectif de sous-groupes dont la limite projective $H$ s'identifie à un sous-groupe de $G$. Il est clair en outre que $H\supset K$ et on voit aussitôt que $K$ est dense dans $H$ (cf. prop.23) ; mais en général on aura $H\neq K$ (exerc. voir §3, cor.1 de la prop.18).
\margin{\unsure{2|?}}

\uline{Corollaire} 3 - \uline{Les hypothèses et notations étant celles du cor.1, soit} $N_\alpha$ \uline{le noyau de} $g_\alpha$ ; \uline{alors} $(N_\alpha)$ \uline{est un système projectif de sous-groupes distingués, dont la limite projective est le noyau de} $g=\varprojlim g_\alpha$

La relation $g_\alpha\circ\varphi_{\alpha\beta}=\psi_{\alpha\beta}\circ g_\beta$ pour $\alpha\leq\beta$ montre aussitôt que $(N_\alpha)$ est un système projectif. En outre, si $\varphi_\alpha$ (resp. $\psi_\alpha$) est l'application canonique de $A$ dans $A_\alpha$ (resp. de $B$ dans $B_\alpha$), la relation $\psi_\alpha(g(x))=g_\alpha(\varphi_\alpha(x))$ pour $x\in A$ montre que si $x\in N=\varprojlim N_\alpha$, on a $\psi_\alpha(g(x))=e$ pour tout $\alpha\in I$, donc $g(x)=e$, et réciproquement cette dernière relation entraîne $\varphi_\alpha(x)\in N_\alpha$ pour tout $\alpha\in I$, donc $x\in N$. Ce raisonnement montre en outre que $N$ est \uline{l'intersection des sous-groupes} $\varphi_\alpha^{-1}(N_\alpha)$ (qui forment une base de filtre, comme on le vérifie aussitôt).

\margin{101}
